\begin{proposition}
\label{proposition-compute-cohomology-crystal}
Assumptions as in Proposition \ref{proposition-compute-cohomology}
but now assume that $\mathcal{F}$ is a crystal in quasi-coherent modules.
Let $(M, \nabla)$ be the corresponding module with connection over $D$, see
Proposition \ref{proposition-crystals-on-affine}. Then the complex
$$
M \otimes^\wedge_D \Omega^*_D
$$
computes $R\Gamma(\text{Cris}(X/S), \mathcal{F})$.
\end{proposition}

\begin{proof}
We will prove this using the two spectral sequences associated to the
double complex $K^{*, *}$ with terms
$$
K^{a, b} = M \otimes_D^\wedge \Omega^a_{D(b)}
$$
What do we know so far? Well, Lemma \ref{lemma-vanishing}
tells us that each column $K^{a, *}$, $a > 0$ is acyclic.
Proposition \ref{proposition-compute-cohomology} tells us that
the first column $K^{0, *}$ is quasi-isomorphic to
$R\Gamma(\text{Cris}(X/S), \mathcal{F})$.
Hence the first spectral sequence associated to the double complex
shows that there is a canonical quasi-isomorphism of
$R\Gamma(\text{Cris}(X/S), \mathcal{F})$ with
$\text{Tot}(K^{*, *})$.

\medskip\noindent
Next, let's consider the rows $K^{*, b}$. By
Lemma \ref{lemma-structure-Dn}
each of the $b + 1$ maps $D \to D(b)$ presents $D(b)$ as the $p$-adic
completion of a divided power polynomial algebra over $D$.
Hence Lemma \ref{lemma-relative-poincare} shows that the map
$$
M \otimes^\wedge_D\Omega^*_D
\longrightarrow
M \otimes^\wedge_{D(b)} \Omega^*_{D(b)} = K^{*, b}
$$
is a quasi-isomorphism. Note that each of these maps defines the {\it same}
map on cohomology (and even the same map in the derived category) as
the inverse is given by the co-diagonal map $D(b) \to D$ (corresponding
to the multiplication map $P \otimes_A \ldots \otimes_A P \to P$).
Hence if we look at the $E_1$ page of the second spectral sequence
we obtain
$$
E_1^{a, b} = H^a(M \otimes^\wedge_D\Omega^*_D)
$$
with differentials
$$
E_1^{a, 0} \xrightarrow{0}
E_1^{a, 1} \xrightarrow{1}
E_1^{a, 2} \xrightarrow{0}
E_1^{a, 3} \xrightarrow{1} \ldots
$$
as each of these is the alternation sum of the given identifications
$H^a(M \otimes^\wedge_D\Omega^*_D) = E_1^{a, 0} = E_1^{a, 1} = \ldots$.
Thus we see that the $E_2$ page is equal $H^a(M \otimes^\wedge_D\Omega^*_D)$
on the first row and zero elsewhere. It follows that the identification
of $M \otimes^\wedge_D\Omega^*_D$ with the first row induces a
quasi-isomorphism of $M \otimes^\wedge_D\Omega^*_D$ with
$\text{Tot}(K^{*, *})$.
\end{proof}